% !PRINTED 594
\weberpage{594}
\begin{center}
\textbf{\S\,160.\qquad \gesperrt{Primideale ersten Grades in $\mathfrak{T}_{\mathfrak{m}}$.}}
\end{center}

Es kommt nun vor allem darauf an, die Primzahlen $p$ aufzusuchen, die im Körper $\mathfrak{L}$ in Primideale ersten Grades zerlegbar sind. Nach § 122 und § 156 sind das die Primzahlen $p$, die durch die Hauptform der Diskriminante $D = 4\Delta$ oder $16\Delta$ darstellbar sind. Diese sind von der Form:
\[
\begin{aligned}
1) \quad & p = y^2 - \Delta x^2, \qquad \Delta \equiv 0 \ (\mathrm{mod}\ 4) \text{ oder } \equiv 1 \ (\mathrm{mod}\ 8),\\
2) \quad & p = y^2 - 4\Delta x^2, \qquad \Delta \equiv 5 \ (\mathrm{mod}\ 8).
\end{aligned}
\]
Ist $\Delta \equiv 1 \ (\mathrm{mod}\ 8)$, so muß der dritte Koeffizient der Form $(a, b, c)$ gerade sein. Ist $\Delta \equiv 0 \ (\mathrm{mod}\ 4)$, so können wir $c$ gerade annehmen, indem wir nötigenfalls zu der Parallelform $(a, b+2a, a+b+c)$ übergehen. Zerfällt also $p$ in die beiden Primfaktoren $\pi, \pi'$, so ist
\weq{1}{\begin{aligned}
1) \quad & \pi = y + x\sqrt{\Delta},\\
2) \quad & \pi = y + 2x\sqrt{\Delta},
\end{aligned}}
und um die Multiplikationsformel § 157, (10), (17) auf $\mu = \pi$ anzuwenden, haben wir im letzteren Falle $x$ durch $2x$ zu ersetzen. Wir können also in allen Fällen $cx$ gerade und daher $\varepsilon = \pm 1$ annehmen.

Die in den Darstellungen (1), 1), 2) liegenden Bedingungen für den Modul 2 lassen sich in die eine zusammenfassen:
\[
\pi \equiv 1 \ (\mathrm{mod}\ 2).
\]
Setzen wir, wenn $N(\pi) = p$ ist, für den Augenblick zur Abkürzung:
\weq{2}{\begin{aligned}
\varphi(x) &= A_1x + A_3x^3 + \cdots + A_{p-2}x^{p-2} + x^p,\\
\psi(x) &= 1 + A_{p-2}x^2 + \cdots + A_3x^{p-3} + A_1x^{p-1}, % source plus has a small damage tick; no mathematical prime is printed here
\end{aligned}}
so ist nach § 157, (18), (19)
\weq{3}{A_1 = \pm\pi,}
und
\weq{4}{\pm \boldsymbol{s}(\pi u) = \frac{\varphi(\boldsymbol{s})}{\psi(\boldsymbol{s})},}
und daraus ergibt sich durch Differentiation nach $u$:
\weq{5}{\pm \frac{\pi \boldsymbol{s}'(\pi u)}{\boldsymbol{s}'(u)} = \frac{\varphi'(\boldsymbol{s})\psi(\boldsymbol{s}) - \varphi(\boldsymbol{s})\psi'(\boldsymbol{s})}{\psi(\boldsymbol{s})^2}.}
Nach § 155, (5) ist aber
\[
\boldsymbol{s}'(u)^2 = 1 - \tfrac{1}{4}\lambda \boldsymbol{s}^2 + \boldsymbol{s}^4,
\]
%CONT>>
